\documentclass[11pt,a4paper]{article}
% Compile with XeLaTeX (fontspec and xeCJK).
\usepackage{fontspec}
\setmainfont{LibertinusSerif}[Extension=.otf,UprightFont=*-Regular,BoldFont=*-Bold,ItalicFont=*-Italic,BoldItalicFont=*-BoldItalic]
\setsansfont{LibertinusSans}[Extension=.otf,UprightFont=*-Regular,BoldFont=*-Bold,ItalicFont=*-Italic]
\usepackage{xeCJK}
\setCJKmainfont{Hiragino Mincho ProN}
\xeCJKDeclareCharClass{Default}{"00B0, "00B7, "2013, "2014, "2018 -> "201F, "2026}
\usepackage[french]{babel}
\usepackage[a4paper,margin=2.5cm,top=2.3cm]{geometry}
\usepackage{graphicx}
\usepackage{xcolor}
\usepackage{titlesec}
\usepackage{fancyhdr}
\usepackage[hidelinks]{hyperref}
\definecolor{accent}{RGB}{31,73,125}
\definecolor{muted}{RGB}{93,104,119}
\hypersetup{colorlinks=true,urlcolor=accent,linkcolor=accent,pdfauthor={Romain Fernex},pdftitle={La femme japonaise dans le Japon de l’ère Meiji : des trajectoires féministes à l’épreuve du conservatisme}}
\setcounter{secnumdepth}{0}
\titleformat{\section}{\large\bfseries\color{accent}}{}{0em}{}
\setlength{\parindent}{1.2em}
\setlength{\parskip}{0.35em}
\pagestyle{fancy}
\fancyhf{}
\fancyhead[L]{\small\color{muted}Romain Fernex}
\fancyhead[R]{\small\color{muted}\itshape La femme japonaise dans le Japon de l’ère Meiji}
\fancyfoot[C]{\small\color{muted}\thepage}
\renewcommand{\headrulewidth}{0pt}
\newcommand{\partbreak}{\par\begin{center}\color{muted}$\ast$\quad$\ast$\quad$\ast$\end{center}\par}

\newcommand{\articlefigure}[2]{
  \begin{center}
  \includegraphics[width=0.55\linewidth,height=6.5cm,keepaspectratio]{#1}
  \par\smallskip{\small\itshape\color{muted}#2\par}
  \end{center}}

\begin{document}
\thispagestyle{empty}
\noindent{\small\sffamily\color{accent}\MakeUppercase{Histoire du Japon}}\par\medskip
\noindent{\LARGE\bfseries La femme japonaise dans le Japon de l’ère Meiji\par}
\smallskip\noindent{\large\itshape Des trajectoires féministes à l’épreuve du conservatisme\par}
\medskip
\noindent{\color{accent}\rule{\linewidth}{0.8pt}}\par\smallskip
\noindent{\small\color{muted}Romain Fernex · 2023 · Article en français}\par\bigskip

Le 8 juillet 1853, le commodore de l’armée de Etats-Unis Matthew Perry, pénètre dans le port de Tokyo abord de la frégate USS Susquehanna afin d’exiger la réouverture du Japon au commerce international. Seulement 15 ans plus tard, la fin de la politique japonaise de Sakoku (politique d’ouverture très restreinte sur le monde extérieur) est officiellement levée et la mort de l’empereur Komei marque le passage de l’ère Tokugawa à l’ère Meiji. Cette dernière se distingue par des changements majeurs sur tous les plans et plus particulièrement sur le plan industriel et politique avec la fin du shogunat. Cependant, au regard du monde, cette révolution semble avoir épargné une partie de la population japonaise, les femmes. Peu importe leur statut, elles apparaissent comme délaissées par ce changement, comme condamnées à rester prisonnières d’une ère dominée par un système de filiation patrilinéaire et des valeurs confucéennes qui limitent drastiquement leurs libertés tant dans la cellule familiale que dans la sphère publique. Est-ce donc à dire que l’ère Meiji n’a rien changé pour les femmes au Japon~? Les progrès de la sociologie et des recherches historiques dans le domaine nous permettent aujourd’hui de dire que non, bien au contraire. Longtemps invisibilisée, la lutte des femmes japonaises face à la à une société japonaise très patriarcale n’a pas attendu les années 1960 ou même l’ère Taisho pour se structurer et penser sa structure idéologique, ses revendications. De fait, dès l’ère Tokugawa, on peut déjà observer des traces de l’engagement de certaines femmes dans la sphère politique même si celles-ci tendent à rester isolées. Mais alors, pourquoi choisir l’ère Meiji~? En quoi celle-ci marque-t-elle un nouveau départ pour la lutte des femmes japonaises vers leur indépendance~? C’est pour répondre à cette question que nous tenterons d’analyser, au cours de cet article, la naissance et les trajectoires, non pas du féminisme mais des féminismes, qui ont pu émerger à cette époque et en présenterons les acteurs majeurs. Féminisme que nous entendrons dans la suite de cet article comme étant «~la prise de conscience des contraintes imposées aux femmes en raison de leur sexe et les tentatives visant à supprimer ces contraintes et à mettre en place un système de genre plus équitable impliquant de nouveaux rôles pour les femmes et de nouvelles relations entre les femmes et les hommes~»{[}1{]}. Dès lors nous considérerons comme féministe tout individu, homme ou femme, qui a conscience de cette inégalité et qui s’exprime contre celle-ci. Enfin, il est important de noter que si le présent usage du terme “féminisme” ne constitue pas un anachronisme à proprement parler, il n’est que très rarement (voir jamais) utilisé par les militantes présentées au cours de cet article. Cela est dû au fait que ce concept est encore récent à l’époque et que son évocation constitue une forte prise de risque dans un cadre conservateur, au Japon comme ailleurs.

Pour répondre à la problématique de cet article, nous suivrons l’ordre chronologique en parlant en premier lieu de la transition entre l’ère Tokugawa et l’ère Meiji avec la redéfinition du rôle des femmes dans un contexte d’agitation politique et de dynamisme intellectuel dû en partie à une ouverture nouvelle sur l’Occident. Puis nous verrons comment ces nouveaux mouvements féministes réagissent à la réaffirmation d’un état conservateur, soucieux de garder la mainmise sur une société en transformation, dans les années 1884 à 1910. Enfin nous nous pencherons sur la fin de l’ère Meiji et l’aboutissement puis la passation d’un corpus idéologique féministe qui survivra à ces premiers mouvements pour venir enrichir la lutte féministe dès ères qui suivent.

\partbreak

Tout d’abord, il convient de faire un état de lieu de la situation de la femme japonaise durant l’ère Tokugawa~: le régime politique inflexible et autoritaire de l’ère Tokugawa se démarque par une mainmise particulièrement forte de l’Etat sur la sphère privée comme publique. La femme japonaise, et notamment les femmes de samurai qui occupent pourtant les hautes sphères de la société japonaise, est réduite à sa fonction procréatrice, n’ayant aucun droit ni sur son mari ni sur ses enfants. Les femmes de plus basses extractions sont un peu moins contraintes sur le plan familial mais demeurent exploitées dans les champs ou les usines, payant de leur temps le peu de liberté qu’elles peuvent avoir par rapport à leurs consœurs plus fortunées. La prégnance de valeurs confucéennes particulièrement hostiles à leur émancipation est visible à travers plusieurs textes importants à l’époque tel que l’Onna Daigaku (greater learning for women) qui vise à décrire la femme Tokugawa idéale. Celle-ci ce voit ainsi affublée de nombreux défauts tels que la paresse ou la désobéissance qui doit faire d’elle l’esclave de son mari, seule personne présumée respectable de la maisonnée. De manière similaire, le principe de Danson Johi (respectez l’homme, méprisez la femme) domine de nombreux aspects de la société de l’époque rendant le concept d’égalité des genres on ne peut plus lointain. Est-ce à dire pour autant que ce modèle est sans contestation~? Anne Walthall, historienne américaine, a réussi à identifié la présence de femmes dans la vie politique l’ère Tokugawa à travers, entre autre, la présence de nombreuse signatures sur des pétitions présentées par certains villages à l’encontre du pouvoir impérial. Cela montre bien l’implication des femmes, au côté des hommes, dans la lutte légale face à des autorités répressives. Par ailleurs, plusieurs femmes, à l’instar de Matsuo Taseko femme de chef de village, ont réussi à dépasser au moins en partie les conventions sociales qui pesaient sur elles et à avoir une réelle présence dans la sphère publique. Cependant, comme le remarque Walthall, ces manifestations restent sporadiques et ne donnent jamais lieu à une convergence des luttes qui aurait pour objectif l’amélioration du statut de la femme japonaise dans son expression la plus générale.

Aussi, le choc de la rencontre de certains Japonais et Japonaises avec des femmes occidentales n’en est que plus brutal. A la fin de l’ère Tokugawa, à l’occasion d’une mission diplomatique du gouvernement japonais à la Maison Blanche, la présence et l’attitude de la nièce du président Américain James Buchanan surprend les émissaires nippons. Celle-ci préside la séance et se montre plus que capable de s’imposer à table alors même que cela serait tout à fait impensable pour une femme japonaise au vu du statut que lui impose la société japonaise. Cette rencontre, une parmi tant d’autre, contribue à la prise de conscience, de la part d’un gouvernement japonais en pleine réformation, que le statut de la femme est un des éléments qui sépare radicalement le monde occidental de l’archipel japonais. Or, le Japon de l’époque est encore traumatisé par les traités inégaux que lui ont imposés les Etats-Unis et cherche plus que jamais à dépasser son retard de développement face à l’Occident afin de regagner son indépendance. Aussi se développe un nouveau discours parmi certaines élites japonaises~: puisque l’éducation et le statut des femmes semblent être des éléments importants de la réussite occidentale, il faut que le Japon se transforme aussi sur ce point. Cette révélation donne lieu à une politique particulièrement ambiguë qui oscille entre mesure ostentatoire pour plaire à l’occident, discours progressistes et relents de conservatisme face à ce qui est perçu comme une remise en question du modèle social traditionnel japonais.

Un des premiers piliers de cette réforme est l’éducation. Déjà en 1872, l’américain David Murray, conseillant le gouvernement japonais, appelait à inclure les femmes dans la réforme du système d’éducation et appelait à la création d'écoles pour les femmes. Un appel qui n’émane pas uniquement de l’extérieur, comme en témoigne l’engagement de certains membres des élites japonaises de l’époque dans la presse. La Meirokusha (la société des six Meiji), fondée en 1873 et particulièrement influente à l’époque, publie plusieurs articles de son journal, le Meiroku Zasshi, sur le sujet des droits des femmes. Parmi ses rédacteurs, plusieurs se démarquent pour leur position critique sur le sujet du traitement des femmes dans la société japonaise, qui passe pour un réel engagement progressiste pour l’époque. Ainsi, le vicomte Mori Arinori, membre fondateur et premier ambassadeur du Japon aux Etats-Unis fait valoir le droit des femmes à l’éducation en prônant le principe de Ryosai Kenbo (bonne épouse, sage mère). Si il est évident que celui-ci se garde bien d’appeler à l’égalité homme-femme, cette affirmation que la femme a un réel rôle dans la famille en tant qu’éducatrice, et qu’elle doit donc être formée pour assumer ce rôle, est particulièrement novatrice pour l’époque. Après un premier mouvement de création d’écoles pour les filles de samurai dans les années 1850, une politique d’éducation ambitieuse est mise en place par le gouvernement en 1872 avec la publication du Gakusei (code fondamental sur l’éducation) qui prône une harmonisation des sujets étudiés en classe peu importe le genre ou la classe sociale et qui décrète l’éducation comme étant un droit pour tous et toutes. Ce mouvement accompagne la création de nombreuses écoles privées réservées aux filles, la Ueda Jogakko à Tokyo en 1872 ou la Tokyo Joshi Shihan Gakko (l’école normale féminine de Tokyo) en 1874 avec le soutien de l’impératrice. Pour autant, celui-ci a aussi ses limites~: si l’école est désormais un droit pour tous, peu nombreuses sont les familles prêtes à financer des études à leurs filles, d’autant plus que l’éducation des filles est encore vue comme négative par une bonne partie de la société. De plus, les filles qui s’inscrivent dans une de ces écoles sont souvent contraintes de les quitter avant d’en être diplômées pour des raisons familiales ou pécuniaires. Toutefois, il est important de mentionner que le gouvernement n’est pas le seul acteur de ce renouveau dans l’éducation des filles. Les Occidentaux jouent aussi un rôle important avec la mise en place d’un réseau d’écoles de missionnaires qui viennent enseigner aux jeunes Japonais et Japonaises les bases de la religion chrétienne{[}1{]} mais aussi de la physique, des mathématiques ou encore de la littérature occidentale. La Ferris Jogakuin (école Ferris pour les filles) établie par la missionnaire Clara Hepburn en 1859 puis reprise par Mary Kidder en est une parmi de nombreuses autres. Celles-ci constituent un large réseau d'institutions ouvertes aux Japonaises peu importe leur classe sociale. Ce réseau, ainsi que celui qui se met en place dans certaines écoles gouvernementales, est particulièrement important pour le développement d’une véritable conscience féministe. De fait, de jeunes femmes se retrouvent entre elles, souvent loin de leur environnement familial, et sont exposées à des ouvrages occidentaux souvent plus ouverts sur les questions liées à la place des femmes{[}2{]}. Cet environnement nouveau, marqué par le temps de l’éducation souvent très formateur pour la jeunesse, aboutit à l’expression d’un nouveau mode de socialisation pour les japonaises et, pour certaines d’entre elles, l’entrée dans la sphère politique en y introduisant des problématiques pensées par les femmes et qui les touchent directement.

\articlefigure{media/mary_kidder-1.jpg}{Mary Eddy Kidder (1834–1910)}

Cette nouveauté politique et sociale s’insère dans un cadre de revendication sociale et de politisation des luttes face au gouvernement. En effet, en 1877 la dernière poche de résistance militaire face au gouvernement Meiji est définitivement éteinte, ne reste alors que la voie politique aux réformateurs. De fait, dès 1873, plusieurs politiciens dont Itagaki Taisuke, ancien samurai, décide de se retirer du gouvernement pour fonder un nouveau mouvement réformateur, le mouvement pour les libertés et les droits du peuple (Jiyu Minken Undo), qui vise à l’établissement d’une véritable démocratie. Ce mouvement, d’abord hostile à l’intégration de femmes en son sein, prend progressivement conscience leur importance en tant que force politique à part entière et plusieurs de ses membres se démarquent par leur position particulièrement ouverte sur les sujets liés aux droits des femmes. Plusieurs des leaders du mouvement inspiré par l’Occident, tel que Ueki Emori, voient dans l’amélioration du statut de la femme, certes principalement au sein du foyer, un réel moteur pour le développement du Japon. Aussi, ce mouvement ouvre-t-il au fur et à mesure des tribunes pour des personnages féminins, y compris et surtout sur des questions liées au statut de la femme. Ainsi en 1881, Kishida Toshiko, jeune japonaise, donne son premier discours en public qui connaît un franc succès. Riche d’une éducation d’excellence, ayant été tutrice de l’impératrice pendant plusieurs années, elle s’élève contre le système de concubinage {[}3{]} qui domine à la cour comme dans les foyers japonais et revendique l’affirmation de la femme dans la sphère tant privée que publique. Venant de Kyoto, elle est particulièrement familière avec le statut précaire des jeunes filles japonaises {[}4{]} et sa connaissance des grandes figures du féminisme occidental, comme Millicent Fawcett et le mouvement des suffragettes, lui permet d’avoir un écho tout particulier auprès de son auditoire. Aussi en octobre 1883 elle organise la première rencontre de la Joshi Daienzetsukai (la société de lecture des femmes kyotoïtes) qui attire près de 2000 femmes. Elle multiplie également les publications pour critiquer le gouvernement et défendre le statut de la femme japonaise en faisant tantôt appel à de grandes figures féminines de l’histoire japonaise comme Murasaki Shikibu, autrice de Genji Monogatari{[}5{]} (Conte du Genji), tantôt à des femmes occidentales comme Elizabeth Browning. Néanmoins, cette multiplication des groupes de réflexion féminins et des magazines dédiés à la question des droits des femmes se heurte rapidement à un gouvernement autoritaire, inquiet face à la montée en puissance d’un discours qui lui est hostile et qui est vu comme une menace envers le modèle traditionnel japonais. La montée des répressions et des divisions provoque l’éclatement du Parti Libéral, premier parti de l’ère Meiji et né du Jiyu Minken Undo, en 1884, assénant un coup sévère aux jeunes mouvements féministes qui, pour la plupart, en étaient très proches et y avaient trouvé un réel moyen d’expression politique.

\articlefigure{media/nakajima_kishida_toshiko-1.jpg}{Toshiko Kishida (1863–1901)}

1884 est en effet une année majeure à bien des égards pour la société japonaise et particulièrement pour les mouvements pour la libération de la femme. De fait, en dépit du durcissement de la censure et des conditions économiques et politiques, cette une année de maturation et d’expansion des mouvements féministes au Japon. Mary Leavitt, missionnaire et représentante de la WCTU (woman’s christian temperance union), arrive au Japon à cette date et tient son premier discours en public à Tokyo seulement 2 ans plus tard. Lors de cet évènement sponsorisé par le Jogaku Zasshi {[}6{]}, un des principaux journaux de l’époque attaché à la condition des femmes, elle s’exprime devant un auditorat composé exclusivement de femmes et défend le rôle de la femme tout en dénoncant les vices qui nuisent à son développement. Elle condamne notamment fermement l’usage de l’alcool et du tabac par les hommes japonais et critique le système de concubinage toujours en vigueur au Japon, malgré plusieurs mesures d’apparat. Si le sous-texte religieux qui sous-tend ses propos n'échappe pas à une bonne partie de l’auditoire, beaucoup de personnes retrouvent en ces mots bien des problèmes qu’elles rencontrent au quotidien. Mary Leavitt promeut par ailleurs l’éducation des femmes en donnant des exemples très concrets de femmes européennes qui commencent tout juste à accéder à des postes jusqu’ici réservés aux hommes comme celui de docteure ou de juriste. Un discours qui se fait dans un contexte de prise de conscience, de la part de certaines féministes japonaises comme Toshiko Kishida, des conditions en réalité difficiles des femmes outre-atlantique. Aussi, loin d’y trouver un motif de désespoir, Kishida y voit plutôt, dans un article publié en 1884, une raison de considerer la lutte féministe comme un moyen de définitivement rattraper l’écart avec l’Occident en réalisant ce que celui-ci n’a pas encore réussi à accomplir. C’est aussi à cette époque que des jeunes femmes inspirées par les discours des missionaires et des femmes comme Kishida se lancent à leur tour dans la lutte pour le droit des femmes, poussées tantôt par le christianisme, tantôt par le socialisme (ou les deux).

Fukuda Hideko est l’une d’entre elles. Encore adolescente à l’époque du discours tenu par Kishida en 1882, elle en retire un prochain engagement pour la cause qu’elle conservera toute sa vie. Après plusieurs tentatives, dont l’établissement d’une école privée pour filles, fermée en 1884 faute de moyens, elle est arrêté par la police en novembre 1885 lors de heurts à Osaka{[}7{]}. Alors qu’elle est en prison, en 1886, Yajima Kajiko, autre militante féministe inspiré par les mots de Mary Leavitt la même année, fonde la Tokyo Fujin Kyofukai (la société de réforme des femmes tokyoïte) qui devient rapidement un organisme important dans la diffusion d’idées féministes au Japon. Peu après sa sortie de prison, en 1889, le gouvernement Meiji vote une première constitution qui se trouve être particulièrement défavorable au statut des femmes. Celle-ci nie aux femmes, à travers l’Article 5 du règlement de sécurité de la Police, tout droit de participer en politique, à l’échelle locale comme nationale. Celui-ci interdit entre autres aux femmes d’assister aux futures réunions du Parlement pour décider des prochaines lois. Cette décision provoque une véritable levée de boucliers dans la société japonaise, et plus particulièrement chez les militantes de l’ancien Jiyu Minken Undo dont sont membres certains des membres du nouveau gouvernement. Yajima Kajiko et la Tokyo Fujin Kyofukai lancent ainsi de nombreuses protestations pour contester l’interdiction aux femmes de pénétrer l’enceinte de la Diète. Cette opération représente une des premières campagnes organisées à l’échelle nationale par un mouvement pour le droit des femmes et la pétition ainsi portée permet un très léger adoucissement de la mesure initiale. Alors que la majorité des activités de la Tokyo Fujin Kyofukai se tournait d’ordinaire vers la condamnation du système de concubinage et des réseaux de prostitution{[}7{]}, il s’agit ici pour les femmes japonaises de penser leur rôle en tant qu’acteur politique dans la sphère publique et non plus seulement dans la sphère privée qui est celle du foyer. C’est dans ce contexte particulièrement agité que Fukuda Hideko fonde une nouvelle école privée pour jeunes femmes et se rapproche du jeune mouvement socialiste dirigé par Sakai Toshihiko. Ce mouvement hétérogène se forme en 1901 en réaction face à la montée de l’autoritarisme du gouvernement japonais, marquée par la réforme du Code civil de 1898 qui se révèle désastreuse pour les femmes du fait de la réaffirmation de la vision réductrice de la femme, qui faisait autorité à l’époque Tokugawa. Cette réforme, qui apparait après la guerre entre la Chine et le Japon en 1895, montre que l’autoritarisme du gouvernement Meiji ne cesse de croître sous la pression, et le retour aux valeurs traditionnelles intervient comme une volonté de revenir à la stabilité. Aussi, le mouvement socialiste regroupe des chrétiens, des anciens du Jiyu Minken Undo, des militantes et autres figures qui ont en commun de voir en le socialisme un moyen d’atteindre le développement désiré par le Japon, et surtout qui sont très critiques de la tournure que prend le gouvernement.

\articlefigure{media/hideko_fukuda_cropped.jpg}{Fukuda Hideko (1865–1927)}

Ce mouvement fait par ailleurs écho à une autre transformation majeure qui débute un peu avant les années 1880 au Japon, l’industrialisation massive du pays. De fait, pour réaliser la modernisation de son économie, le Japon s’est tourné vers la production et l’exportation de soie, ouvrant de multiples moulins et usines dédiés à cet effet. Or de 1894 à 1912, 60\% de la main-d'œuvre employée dans l’industrie au Japon est féminine. En effet, nombre de jeunes femmes de familles pauvres sont enrôlées de force dans les exploitations du pays où elles sont confrontées à des conditions de travail particulièrement difficiles, comme en témoigne le film Nomugi Pass de Satsuo Yamamoto, paru en 1979. Les ouvrières ne reçoivent qu’une maigre compensation financière qu’elles donnent immédiatement à leur famille, ne leur permettant pas d’acquérir le moindre soupçon d’indépendance financière en dépit de leurs efforts. Cette difficile réalité s’accompagne néanmoins d’importants mouvements de contestations à l’échelle locale, marqués notamment par d’importantes grèves dans certaines exploitations. En juin 1886 à Amamiya, des ouvrières, bientôt rejointes par leurs quelques collègues masculins travaillant dans les moulins, s’élèvent contre la baisse des salaires, et en faveur d’une hausse des primes, qui constituent souvent leur seule source réelle de revenu. Cette mobilisation, la première d’une telle ampleur à l’ère Meiji, dure plusieurs semaines et aboutit à une victoire de la part des grévistes. Si celle-ci est la première d’une longue série, toutes ne rencontrent pas le succès et aucun mouvement national ne parvient à prendre forme à cause du système de dortoirs qui isolent les ouvrières. De plus, la différence de milieu social entre ces ouvrières et la majorité des militantes féministes, qui pour la plupart sont issues de la classe moyenne et ont reçu un minimum d’éducation, limite fortement la communication entre ces deux corps.

Mais alors qu’en est-il du mouvement féministe socialiste, entre rupture et continuité avec les nombreux mouvements de femmes de classes moyenne supérieure, souvent teinté de christianisme, comme celui porté par Toshiko Kishida~? Plusieurs membres du mouvement ont déjà été impliqués dans la lutte sociale en poussant à la création de syndicats en 1897 face à la crise économique. Pour ce qui est des femmes qui ont rejoint le mouvement, le besoin d’imposer davantage d’enjeux liés à la condition des femmes les poussent à former un groupe à part entière pour se concentrer sur ces problèmes. Elles sont aidées dans cet effort par certains leaders du mouvement qui imposent qu’à chaque meeting tout homme soit accompagné d’une femme afin de préserver une forme de parité, et permettent à plusieurs femmes de prendre la parole publiquement. Cependant, l’association est dissoute en 1904 face aux pressions du gouvernement qui se font de plus en plus insistantes dans un contexte de guerre avec la Russie et de division internes entres chrétiens, matérialistes et autres tendances au sein du mouvement. Cependant la fin du féminisme socialiste intervient réellement à la suite de “l’incident des drapeux rouges” le 22 juin 1908. La libération, par le gouvernement, de l’activiste et militant anarchiste Koken Yamaguchi entraîne une importante mobilisation du mouvement socialiste qui accueille la nouvelle avec joie en organisant une manifestation non autorisée dans les rues de Tokyo. Parmi les membres du mouvement, plusieurs sont de fervents soutiens des idées de Yamaguchi, dont Kanno Suga, et brandissent des drapeaux rouges avec des slogans anarchistes. Cette mobilisation entraine une répression intense de la part de la police. Le magazine du mouvement est interdit et Kanno Suga est arrêtée à peine un an plus tard pour avoir plannifié l’assassinat de l’empereur. Celle-ci est exécutée en 1910 avec 24 autres membres du courant marquant définitivement la fin de l’aventure socialiste.

\partbreak

Cette vague de répression a mis fin à de nombreux mouvements féministes en rendant leurs activités très risquées. Cependant, nombre de ces militantes n’ont pas abandonné la lutte et, bien que le contrôle sur les écoles missionnaires et sur la presse soit plus intense, les publications sur le sujet sont encore relativement régulières, même si elles restent bien souvent à l’échelle locale et sont l’objet de censure. En 1911, malgré la vague de répression contre les mouvements d’opposition, est fondée l’association Seito (plus connue sous le nom de Bluestockings). Celle-ci est unique en ce qu’elle n’a pas à l’origine de motivation politique. En effet, sa fondatrice Hiratsuka Raicho, fille d’une riche famille libérale Tokyoïte et éduquée à l’université féminine de Naruse Jinzo, souhaitait avant tout en faire une tribune pour exposer le travail de femmes artistes et n’est, du moins à l’origine, que peu sensibilisée aux problèmes liés au statut de ses consœurs. De fait, être une écrivaine ou une artiste à succès demeure particulièrement compliqué hormis quelques rares exceptions, et les tribunes offertes aux femmes en matière d’art comme de politique sont particulièrement limitées. Dès sa première édition, le magazine reçoit d'importants soutiens de la part d’artistes établis, et de la part de la mère de la fondatrice qui l’aide financièrement. Naturellement, l’exaltation du talent féminin, peu importe le domaine, ne manque pas de susciter l’ire de plusieurs commentateurs qui dénigrent cette création. Néanmoins, le magazine rencontre beaucoup de succès auprès de jeunes femmes talentueuses, mais aussi venant de diverses classes sociales. Parmi les nouvelles têtes qui rejoignent l’association se trouve Ito Noe, jeune japonaise de basse extraction venue étudier à Tokyo après s’être enfuie de sa région natale de Kyushu suite à un mariage arrangé. Cette dernière milite pour que le journal prenne une tournure davantage politique quitte à risquer la confrontation avec le gouvernement. En avril 1912, un numéro du journal portant un regard critique sur le système familial japonais lui vaut d’être banni temporairement par le gouvernement qui le met désormais sous surveillance de manière plus ou moins officieuse. Cependant, cela n'empêche pas le journal de prendre une tournure toujours plus politique avec une série complète sur les problèmes des femmes un an plus tard. Cette série rencontre un véritable succès auprès de son lectorat et son contenu est relayé dans plusieurs journaux nationaux d’envergure. Parmi les thématiques abordées par le journal, on peut citer le mouvement féministe international à travers le mouvement de suffragettes, la lutte pour l’autonomie et l’égalité des droits ou encore la sexualité des femmes. Fukuda Hideko entre d’ailleurs en contact avec l’association, prônant le socialisme comme remède à ces problèmes, ce qui lui vaudra des réactions mitigées de la part de certains de ses membres.

\articlefigure{media/the_first_issue_of_seito_-_september_1911.jpg}{Couverture de la première édition de Seitō, septembre 1911}

La situation se tend encore plus quand le Ministère de l'Intérieur promulgue une liste de magazines féminins interdits de publication plus tard dans l’année. Si les Bluestockings n’en font pas partie, l’association décide de renforcer encore un peu plus son engagement pour la cause en modifiant ses statuts. Cependant, cette direction toujours plus politique inquiète plusieurs des membres qui ont rejoint l’association principalement pour que leur travail gagne en visibilité. Ces dissensions ne cessent d’augmenter jusqu’en 1915, date à laquelle Raicho, épuisée, cède les commandes de l’association à Ito Noe. Isolée et disposant de ressources financières très limitées, Ito Noe parvient tout de même à faire fonctionner l’association pour un an encore. Une année qui se révèle particulièrement prolifique puisqu’elle parvient à introduire plusieurs questions importantes dans le débat publique comme celle de l’avortement ou de la parentalité. Des questions qui se posent encore, bien après la fin du journal. La fin est tragique pour Ito Noe, arrêtée par la police lors du tremblement de terre de Tokyo en septembre 1923, elle est tuée en cellule avec son neveu et son mari par un policier radicalisé. Néanmoins, les bluestockings laissent derrière elle un héritage puissant, tirant certes parfois inconsciemment, des décennies de luttes qui les ont précédés et posant les bases de la seconde vague de féminisme qui touche l’ère Taisho (qui commence officiellement en 1912). Ainsi en janvier 1920, Raicho, accompagnée de Ichikawa Fusae et Oku Mumeo crée la Shinfujinkyokai (la nouvelle société des femmes) ainsi que son magazine associé, Josei Domei (la ligue des femmes). Cette société, riche de l’expérience des femmes de l’ère Meiji et du début de l’ère Taisho, est un des acteurs majeurs de la lutte féministe pour les décennies qui suivent.

\partbreak

En conclusion, loin d’être une période de stagnation et de renoncement pour les femmes, l’ère Meiji est marquée par l’apparition de nouvelles formes de socialisation féminines et par l’interconnexion de lutte féministes jusqu’à présent désorganisées. Si le climat politique n’a de toute évidence globalement pas été favorable à la moindre réforme dans le sens de l’égalité des genres, cette expérience se révéle précieuse pour les mouvements à venir, et de nombreuses bases de mouvements féministes actuels au Japon y trouvent des fondations idéologiques solides. Ainsi, les apports occidentaux mélangés à une «~conscience de classe~» émergente et aidée par la multiplication des médias, grâce à la démocratisation des journaux, ont permis de former une génération de mouvements féministes, certes très divers et parfois critiques les uns envers les autres mais tous unis par une expérience commune de ce que c’est qu’être une femme au Japon, et déterminés à redonner à la femme la place qu’elle mérite.

\textbf{Notes~:}

\textbf{{[}1{]}} Je reprends ici la définition proposée par la chercheuse Margot Badran dans son ouvrage \emph{Feminists, islam and nation~: Gender and the making of modern Egypt}. Celle-ci est par ailleurs reprise dans l’ouvrage de Mara Patessio cité en bibliographie.

\textbf{{[}2{]}} Il ne s’agit bien sûr pas de dire que la situation de la femme est radicalement meilleure en Occident ou que les textes produits en Occident à l’époque sont des exemples de féminisme. Cependant, ceux-ci apportent bien souvent une perspective radicalement différente à des jeunes japonaises jusqu’à présent exposées majoritairement à des textes d’influence confucéenne, qui se veulent avant tout prescripteurs d’une morale qui les excluent de toute velléité de penser la société. A ce titre, des ouvrages comme \emph{Social Statics} du philosophe anglais Herbert Spencer ou \emph{The Subjection of Women} de John Stuart Mill, publiés respectivement en 1851 et en 1869, constituent de véritables révélations pour certaines japonaises qui s’engagent pour la cause féministe.

\textbf{{[}3{]}} La polygamie était un motif récurrent dans les foyers japonais pendant l’ère Tokugawa et pendant l’ère Meiji. Celle-ci se caractérise par la présence de concubines qui sont souvent traitées de la même façon que l’épouse du maître de maison et sont parfois amenées à donner un héritier à ce dernier. Tel héritier est généralement considéré comme légitime et la femme mariée se retrouve souvent à devoir s’en occuper à la demande de son mari. Cette pratique contribue, dans le cadre de la société japonaise, à l’affaiblissement du statut de l’épouse et à l’asymétrie qui pèse dans la vie de couple au Japon. De fait, la lutte contre le concubinage est un thème récurrent de la lutte féministe au Japon lorsque ceci à touche à la sphère privée. Cependant, la forme de cette lutte tend à varier d’un groupe à l’autre et d’une période à l’autre. En effet, si l’attaque du système de concubinage a d’abord visé les hommes et les femmes qui en prenaient part, sans distinction de situation, plusieurs féministes, telle Fukuda Hideko, ne manquent pas de développer une critique plus nuancée et surtout qui prend en compte le statut particulièrement précaire de bien des concubines.

\textbf{{[}4{]}} Le fait que Toshiko vienne de Kyoto a d’autant plus d’importance que la région Kyoto Osaka est le théâtre de pratiques particulièrement dures de la part des hommes. Parmi celles-ci on peut citer le phénomène des «~Daughters in boxes~» où de jeunes filles à marier sont enfermées chez elles par leur famille jusqu’au moment où elles doivent être présentées à leurs futurs époux. Cette remise en question radicale de la liberté des femmes sera vivement critiquée par Kishida dans plusieurs de ses discours, dont un de ses plus célèbres le 12 octobre 1883 à Otsu. Ce discours lui vaut d’ailleurs d’être arrêtée par la police et emprisonnée pendant plusieurs jours. 

\textbf{{[}5{]}} Genji Monogatari est une des œuvres de la littérature japonaise les plus célèbres au Japon comme à l’international. Il est écrit au 11\textsuperscript{ème} siècle, à l’ère Heian, par Murasaki Shikibu, dignitaire de la cour impériale et une des premières poètes à écrire en langue japonaise. De fait, la langue chinoise a longtemps dominé la cour impériale japonaise au détriment du japonais qui était relégué au second rang. Ainsi, une bonne partie des premiers auteurs à succès en langue japonaise sont des femmes de la cour, ce qui témoigne de l’érudition de celles-ci et d’un sens de l’esthétisme qui rivalise tout à fait avec celui de leurs homologues masculins.

\textbf{{[}6{]}} Le Jogaku Zasshi est un des journaux progressistes à grand tirage de l’époque Meiji. Sa publication commence en juillet 1885, sous l’impulsion d’un homme, Iwamoto Yoshiharu, défenseur de l’éducation pour les femmes et instructeur à la Meiji’s Girls school. Nombre d’articles favorables à la cause féministe y sont publiés, dont une majorité écrits par des femmes engagées dans la lutte. Son influence dure jusqu’en 1904, date de sa fermeture.

\textbf{{[}7{]}} La question des réseaux de prostitution est, avec le système de concubinage, un des thèmes récurrent de la lutte féministe au Japon. Celle-ci est particulièrement vaste et mériterait à n’en pas douter un article à part entière, mais voici tout de même quelques précisions utiles~: La société de l’ère Tokugawa comme de l’ère Meiji et de l’ère Taisho est marquée par la mise en place d’un véritable réseau de prostitution institutionnalisé qui consiste en l’envoi aux quatres coins du globe de nombreuses jeunes Japonaises disposant de peu de moyens. Celles-ci se voient ainsi contraintes d’offrir leurs services tantôt à des riches clients étrangers, tantôt à des soldats japonais épuisés par les conquêtes menées par le pouvoir impérial. C’est par ailleurs à l’ère Meiji que cette «~industrie~» se structure avec le phénomène des Karayuki et Ameyuki-san qui se développe fin 19\textsuperscript{ème} et qui persistera pendant la première moitié du 20\textsuperscript{ème} siècle. Ce sont ces structures qui pose les bases du très controversé système des “femmes de réconfort”, mis en place par l’empire japonais dans ses territoires conquis en Asie pendant la seconde guerre mondiale.

\textbf{Bibliographie~:}

\begin{itemize}
\item
  Sharon L. Sievers, \emph{Flowers in Salt, The beginnings of Feminist consciousness in Modern Japan}, Stanford University Press, 1983
\item
  Mara Patessio, \emph{Women and Public life in Early Meiji Japan, The development of the feminist movement}, vol.71, University of Michigan Center for Japanese Studies, 2011
\item
  Anne Walthall, “Devoted Wives/Unruly Women: Invisible Presence in the History of Japanese Social Protest,” Signs 20.1 (1994): 106–36; “Edo Riots,” dans \emph{Edo and Paris: Urban Life and the State in the Early Modern Era}, édité par James L. McClain, John M. Merriman, et Ugawa Kaoru, 407–28 (Ithaca: Cornell University Press, 1994).
\item
  Margot Badran, \emph{Feminists, Islam and Nation~: Gender and the Making of Modern Egypt}, PUP, 1995, 19-20
\end{itemize}

\end{document}
